\chapter{Blocks and Closures}
\label{ch:blocks}

\standfirst{A block is a piece of code you can pass around and run later. It
is also the only reason a language with no control structures can express
one.}

\section{The basics}

Square brackets make an object out of a piece of code. The code does not run
until you ask.

\begin{st}
[3 + 4]                  "a BlockClosure; nothing has been added"
[3 + 4] value            "7"
\end{st}

Arguments are named after a colon at the start, followed by a vertical bar:

\begin{st}
[:x | x * x] value: 9                    "81"
[:a :b | a + b] value: 1 value: 2        "3"
[:a :b :c | a + b + c] value: 1 value: 2 value: 3
\end{st}

There are \ct{value}, \ct{value:}, \ct{value:value:},
\ct{value:value:value:} and \ct{value:value:value:value:}. Past four, or when
the count is not known until run time:

\begin{st}
[:x :y | x + y] valueWithArguments: #(3 4)      "7"
\end{st}

A block knows how many arguments it wants, and evaluating it with the wrong
number is an error:

\begin{st}
[:x | x] numArgs          "1"
\end{st}

\section{Block-local temporaries}

A block may declare its own temporaries, after the arguments. They are fresh
at every activation---nilled each time the block is entered, not shared
between calls.

\begin{st}
[:x |
    | doubled |
    doubled := x * 2.
    doubled + 1] value: 3                "7"
\end{st}

This is a post-Blue-Book extension and it works here. A block with no
arguments but with temporaries starts straight in with the bar:

\begin{st}
[ | t | t := 10. t squared ] value       "100"
\end{st}

\section{What a closure closes over}

A block captures the variables of the place it was written, not their values.
It sees later changes, and it can make changes the enclosing method sees.

\begin{st}
| count increment |
count := 0.
increment := [count := count + 1].
increment value.
increment value.
count                                     "2"
\end{st}

The captured variable outlives the method that declared it. This is what makes
them \emph{closures} rather than merely nested code, and it is the single
biggest semantic difference between this system and the Blue Book, where a
block that escaped its method was undefined behaviour.

\begin{st}
"a counter factory: each call answers an independent counter"
makeCounter
    | n |
    n := 0.
    ^[n := n + 1]
\end{st}

Each call to \ct{makeCounter} makes a fresh \ct{n}, so two counters do not
interfere.

\begin{stnote}[Turning closures on]
Full closures are compiled only in the \ct{closures} dialect, which every
profile in \ct{profiles/} except \ct{bluebook} selects---for its packages
and for whatever you hand to \ct{-eval} alike. On the command line,
\ct{-closures} forces them on for a profile that does not ask. If you build
with \ct{-manifest sources/MANIFEST} you get the 1983 dialect, where blocks
are \ct{BlockContext} objects with the 1983 restrictions. Accepting a method
in the Browser, loading a Tonel file into a running server and evaluating in
a workspace all reach the same compiler, so a block means the same thing
wherever it was compiled.
\end{stnote}

\section{Non-local return}

A \ct|^| inside a block returns from the \emph{method the block was written
in}---not from the block. This is how \ct{detect:ifNone:} and every early exit
in the library is written:

\begin{st}
firstNegativeIn: aCollection
    aCollection do: [:each |
        each < 0 ifTrue: [^each]].       "returns from firstNegativeIn:"
    ^nil
\end{st}

The block is being run inside \ct{do:}, several frames down, and the
\ct|^each| unwinds all of that and returns from
\ct{firstNegativeIn:}. Blocks that merely \emph{fall off the end} answer their
last statement, which is why \ct{[3 + 4] value} is 7 with no caret anywhere.

If the method a block belongs to has already returned, a non-local return from
that block has nowhere to go and raises an error---\ct{cannotReturn:}. This is
a real thing that happens when you store a block with a \ct|^| in it and
run it later.

\section{\texttt{ensure:} and \texttt{ifCurtailed:}}

Cleanup that must happen whether or not the protected code finishes:

\begin{st}
[ file writeEverything ] ensure: [ file close ]
\end{st}

\ct{ensure:} runs its argument on the way out no matter how the receiver
leaves: normally, by a non-local return, or by an exception unwinding past it.
\ct{ifCurtailed:} runs its argument \emph{only} on the abnormal exits.

\begin{st}
[3] ensure: [4]                  "3 -- the value is the receiver's"
\end{st}

These two are what make resource handling correct, and they are why
\ct{Mutex} can exist at all: a critical section whose release is in an
\ct{ensure:} cannot be skipped by an exception.

\section{Blocks are objects}

Nothing about a block is special to the compiler once it exists. You can put
one in a collection, store it in an instance variable, pass it to another
process, or send it messages:

\begin{st}
| operations |
operations := Dictionary new.
operations at: #double put: [:x | x * 2].
operations at: #square put: [:x | x * x].
(operations at: #square) value: 7        "49"
\end{st}

Useful protocol beyond \ct{value}:

\begin{center}
\small
\begin{tabular}{@{}ll@{}}
\toprule
\ctp{numArgs} & how many arguments it wants\\
\ctp{whileTrue:}, \ctp{whileFalse:} & loop while the receiver answers so\\
\ctp{repeat} & loop forever; exit with a non-local return\\
\ctp{on:do:} & run under an exception handler\\
\ctp{ensure:}, \ctp{ifCurtailed:} & cleanup\\
\ctp{fork}, \ctp{forkAt:} & run as a separate Smalltalk process\\
\ctp{cull:} & value it, ignoring the argument if it takes none\\
\bottomrule
\end{tabular}
\end{center}

\ct{cull:} deserves a mention because it makes APIs forgiving: a method that
wants to hand your block an argument can use \ct{cull:} and accept both
\ct{[:e | ...]} and \ct{[...]}.

\section{Two things the compiler does behind your back}

Some blocks never become objects at all. When the compiler can see that a
block is used immediately and cannot escape---the arguments of \ct{ifTrue:},
\ct{and:}, \ct{whileTrue:} and friends---it \emph{inlines} the code and emits
a jump. No block is created and no message is sent.

This is why \ct{ifTrue:} costs nothing, and it is also why a rare error
message reads strangely: if the receiver of an inlined \ct{ifTrue:} is not a
Boolean, the failure surfaces as \ct{must be a boolean} from the jump
bytecode rather than as a doesNotUnderstand.

\begin{stwarn}
Because they are inlined, \ct{ifTrue:} and friends require \emph{literal}
blocks to get the fast path. \ct{aBoolean ifTrue: aBlockVariable} still works
--- \ct{True} and \ct{False} really do implement \ct{ifTrue:} --- but it
compiles to a real send.
\end{stwarn}

\section{A frame is only so big}

A method and its blocks share a frame, and the frame has a fixed ceiling:
arguments plus copied values plus temporaries plus working stack must fit. In
practice you meet this only in machine-generated code or a method that has
grown far past the point where it should have been split.

There is a second, firmer ceiling beside it: a method may reference at most
\textbf{63 literals}, because the header holds that count in six bits. Every
distinct string, symbol, large integer and literal array counts, and two
occurrences of \ct{'x'} count twice---they are two mutable objects, and
sharing them would make writing into one change the other. Past 63 the
compiler refuses, naming the method:

\begin{plain}
more than 63 literals referenced; split this method
\end{plain}

The 1983 compiler refuses at the same place, in \ct{Encoder>>litIndex:}, with
the same advice.
